\documentclass[letterpaper]{article}
\usepackage[submission]{aaai2027}
\usepackage[hyphens]{url}
\usepackage{graphicx}
\urlstyle{rm}
\def\UrlFont{\rm}
\usepackage{natbib}
\usepackage{caption}
\frenchspacing
\usepackage{algorithm}
\usepackage{algorithmic}
\usepackage{booktabs}
\usepackage{amsmath}
\usepackage{amssymb}
\pdfinfo{
/TemplateVersion (2027.1)
}
\setcounter{secnumdepth}{0}

\title{Component-Geometry Regularization for Target-Preserving Infrared Small Target Detection}
\author{Anonymous Submission}
\affiliations{}

\begin{document}

\maketitle

\begin{abstract}
Infrared small target detection requires separating tiny, low-contrast objects from clutter while preserving object-level detections and suppressing false alarms. Recent deep models improve dense mask prediction, but their training objectives often supervise the final mask without explicitly regularizing component geometry such as center, boundary, scale, and peak response. This paper drafts CGA-v2, a training-time component-geometry regularizer for MSHNet-style infrared small target detection. CGA-v2 attaches four auxiliary heads to an explicitly declared decoder feature and supervises them with deterministic targets derived from connected foreground components, while inference still uses only the original final mask logit. The current evidence package is intentionally fail-closed: paper evidence is permitted only when dataset preflight, hard-clutter validation source audit, real auxiliary heads, and fixed-threshold final-logit evaluation all pass. In the current repository state, the MSHNetOHEM seed42 epoch400 baseline is verified on NUDT-SIRST, while the paired CGA seed42 epoch400 result, multi-seed statistics, and ablations remain marked as TBD until the matching current-repository runs finish.
\end{abstract}

\section{Introduction}

Single-frame infrared small target detection (IRSTD) aims to recover small, dim target masks from infrared imagery and is important for long-range sensing, surveillance, and search systems where the target may occupy only a few pixels. The task differs from generic semantic segmentation because foreground evidence is sparse, target texture is weak, and clutter can form target-like responses. Classical local-contrast and low-rank formulations exploit this structure through handcrafted priors \citep{chen2014local,dai2017ript,zhang2019pstnn}, while modern neural approaches learn dense predictors and multi-scale representations for the same sparse-object regime \citep{wang2019miss,dai2021acm,dai2021alcnet,li2023dnanet}.

Recent IRSTD models increasingly encode task-specific structure rather than treating the problem as ordinary binary segmentation. Dense nested interaction preserves small objects through repeated high-low feature exchange \citep{li2023dnanet}, shape reconstruction emphasizes edges and target extent \citep{zhang2022isnet}, nested U-shaped feature aggregation improves local-global interaction \citep{wu2023uiunet}, point-supervised label evolution studies how target masks emerge from weak labels \citep{ying2023lesps}, and scale/location-sensitive supervision directly addresses target size and position \citep{liu2024sls}. These methods indicate that where a component is, how large it is, and how its boundary is represented are not secondary details; they are part of the detection signal.

The remaining issue is how to inject such component-level geometry into a strong mask detector without changing the final inference path or making evidence ambiguous. Adding post-processing, auxiliary predictions at test time, or implicit tensor parsing can make an improvement difficult to attribute and hard to reproduce. This concern is especially important for AAAI-style evidence: if the method is a regularizer, the paper should prove that the final mask predictor improves under a fixed protocol, not that extra test-time heads or threshold selection produce better numbers.

We propose CGA-v2, a component-geometry regularization framework for MSHNet-style IRSTD. CGA-v2 keeps the MSHNet final logit as the only inference output and adds four training-only auxiliary heads for center, boundary, scale, and peak supervision. The key implementation constraint is fail-closed: the detector adapter must explicitly declare the final logit and the selected CGA feature, the wrapper must expose all four auxiliary logits, and paper-mode loss construction fails if any required head is absent.

This draft makes three contributions. First, it formulates component geometry as deterministic supervision derived from connected foreground components, producing center, boundary, scale, and peak targets without dataset-specific heuristics. Second, it implements CGA-v2 as a training-only regularizer with explicit output contracts and final-logit-only evaluation, avoiding hidden test-time changes. Third, it defines an evidence-first protocol on NUDT-SIRST that reports both the full test split and a frozen hard-clutter validation split; the current draft includes verified MSHNetOHEM seed42 baseline results and reserves paired CGA claims until current-repository CGA epoch400 evidence is complete.

\section{Related Work}

IRSTD has a long history of methods that separate targets from structured background using local contrast, sparsity, and low-rank assumptions. Local contrast methods amplify target-background differences \citep{chen2014local}, while patch-tensor and nuclear-norm formulations model background redundancy and target sparsity \citep{dai2017ript,zhang2019pstnn}. These approaches make the small-target prior explicit, but they are limited by handcrafted assumptions and often struggle when clutter and targets share local statistics.

Deep IRSTD models replace handcrafted separation with learned dense prediction while retaining task-specific inductive biases. ACM introduces asymmetric contextual modulation for exchanging global context and low-level detail \citep{dai2021acm}; ALCNet combines local contrast with attention \citep{dai2021alcnet}; DNANet develops dense nested interaction and contributes the NUDT-SIRST benchmark \citep{li2023dnanet}. Later models strengthen feature compensation, pyramid context, and transformer-style interaction for complex backgrounds \citep{zhang2022fc3net,zhang2023agpcnet,chen2024tciformer}. CGA-v2 is not positioned as a new backbone family; it regularizes a fixed MSHNet-style detector so that component geometry shapes training while the mask inference path remains unchanged.

Shape-, scale-, and location-aware supervision is the closest line to CGA-v2. ISNet shows that target shape and edge reconstruction matter for infrared small targets \citep{zhang2022isnet}; MSHNet with scale and location sensitivity is the closest architectural and supervision threat because it explicitly targets scale and position in the same task family \citep{liu2024sls}; PConv with SD Loss further studies spatial pixel distribution and scale-dependent loss weighting \citep{yang2025pconv}. CGA-v2 differs by decomposing regularization into four component-derived targets and by making the paper-evidence contract part of the method boundary: auxiliary heads are training-only, and all reported numbers must come from the final logit under a fixed threshold.

Weakly supervised and foundation-model-driven IRSTD further show that target geometry can be separated from dense mask annotation. LESPS studies how point labels can evolve into mask supervision \citep{ying2023lesps}, and SAIST adapts a segmentation foundation model to infrared small targets with language-image guidance \citep{zhang2025saist}. These works broaden the supervision and model-design space. CGA-v2 stays in the fully supervised, single-frame setting and focuses on whether component geometry can improve a controlled MSHNet comparison without changing deployment behavior.

\section{Method}

\subsection{Problem Setting}

Given an infrared image $x \in \mathbb{R}^{1 \times H \times W}$, the detector predicts a binary target mask $y \in \{0,1\}^{H \times W}$. The baseline detector produces a final logit $z \in \mathbb{R}^{1 \times H \times W}$, and evaluation thresholds $\sigma(z)$ at a fixed value of 0.5. CGA-v2 does not add post-processing, connected-component filtering, threshold search, or auxiliary-head inference. Its only change is the training loss applied to a declared intermediate feature.

\begin{figure}[t]
\centering
\fbox{\begin{minipage}{0.92\linewidth}
\centering
\small
Image $\rightarrow$ MSHNet adapter $\rightarrow$ final logit $z$ $\rightarrow$ sigmoid + threshold 0.5 \\
\vspace{0.4em}
Selected decoder feature $\rightarrow$ CGA auxiliary head $\rightarrow$
center, boundary, scale, peak logits \\
\vspace{0.4em}
Training only: geometry losses are added to OHEM + SoftIoU. \\
Inference: auxiliary heads are ignored.
\end{minipage}}
\caption{CGA-v2 regularizes training through component-geometry heads while preserving the final MSHNet mask inference path.}
\label{fig:pipeline}
\end{figure}

\subsection{Fail-Closed Detector Contract}

The detector output is treated as a typed contract rather than an arbitrary tuple of tensors. The MSHNet adapter must expose a dictionary containing the final logit and a single selected CGA feature. In the current implementation, the final logit comes from \texttt{base\_logits/base\_logit}, and the selected feature comes from \texttt{decoder\_features.x\_d0}. The adapter also records metadata for the backbone name, logit source, feature source, feature stride, channel count, and spatial resolution.

This contract prevents silent paper-mode behavior changes. A paper-evidence run cannot select the first four-dimensional tensor as a feature, cannot infer logits from an undocumented output position, and cannot accept a missing auxiliary head in strict CGA mode. The wrapper records the regularizer implementation as \texttt{center\_boundary\_scale\_peak} and explicitly marks \texttt{fallback\_regularizer\_used=false}. The model itself does not declare \texttt{paper\_evidence\_allowed}; that flag is runner/evidence metadata depending on the data and run gates.

\begin{table}[t]
\centering
\caption{Fail-closed method contract used before a CGA result can be treated as paper evidence.}
\label{tab:contract}
\small
\begin{tabular}{lp{0.56\linewidth}}
\toprule
Contract item & Required behavior \\
\midrule
Final logit & Explicit \texttt{logits/final\_logit}; no tuple-position inference. \\
CGA feature & Exactly one declared feature with source, stride, channels, and resolution. \\
Auxiliary heads & All four CGA logits must exist in paper mode. \\
Fallback path & Any fallback or mock regularizer makes the run smoke-only. \\
Evidence flag & Owned by runner metadata, not by the model wrapper. \\
Evaluation & Sigmoid of final logit, threshold 0.5, no auxiliary test-time use. \\
\bottomrule
\end{tabular}
\end{table}

\subsection{Component-Geometry Targets}

CGA-v2 constructs deterministic supervision from the binary mask. Let $\mathcal{C}(y)$ be the set of eight-connected foreground components in a training mask. For each component $C \in \mathcal{C}(y)$, we compute an integer center from the rounded mean of foreground pixel coordinates. The center target $t_c$ marks the center and dilates it by radius $r_c=2$. The boundary target $t_b$ is the binary difference between a radius-$r_b$ dilation and radius-$r_b$ erosion of the foreground mask, with $r_b=1$. The scale target $t_s$ assigns every pixel in component $C$ the value
\begin{equation}
t_s(p) = \min(1, |C| / A_{\max}), \quad p \in C,
\end{equation}
where $A_{\max}=512$ in the default configuration. The peak target $t_p$ is a radius-$r_p=1$ dilation of the center target.

The target builder uses only the current binary mask and fixed radii. It does not depend on dataset names, image identifiers, evaluation splits, or hard-clutter labels. This design keeps the regularizer auditable and separates the method from the evidence protocol.

The four targets are intentionally complementary. The center target supplies a sparse localization cue that is less sensitive to the exact contour of tiny objects. The boundary target pushes the decoder feature to retain the transition between target and background, which is often where false positives and fragmented masks appear. The scale target gives a coarse component-size signal without creating a separate scale classifier, and the peak target encourages a compact high-response region around the component center. None of these maps is used as a prediction at test time; they only shape the feature representation during training.

This decomposition also keeps the ablation logic simple. Removing one target leaves the remaining targets well defined and preserves the same final inference path. Consequently, an ablation can test whether each geometry cue helps the final mask predictor rather than testing a different post-processing rule or a different output head.

\subsection{Auxiliary Heads and Loss}

Given the selected decoder feature $F$, the CGA auxiliary head applies a shared $3 \times 3$ convolution, batch normalization, and ReLU, followed by four $1 \times 1$ prediction heads:
\begin{equation}
(z_c, z_b, z_s, z_p) = h_{\theta}(F).
\end{equation}
The four logits correspond to center, boundary, scale, and peak targets. The baseline loss combines online hard-example-mining binary cross entropy and SoftIoU on the final logit. CGA-v2 adds weighted binary cross entropy losses for the four auxiliary targets:
\begin{equation}
\begin{aligned}
\mathcal{L} ={}& \mathcal{L}_{\mathrm{OHEM}}(z,y)
 + \mathcal{L}_{\mathrm{IoU}}(z,y) \\
& + \alpha(e) \sum_{q \in \{c,b,s,p\}}
 \lambda_q \mathrm{BCE}(z_q,t_q).
\end{aligned}
\end{equation}
where $(\lambda_c,\lambda_b,\lambda_s,\lambda_p)=(0.05,0.03,0.02,0.03)$ in the default protocol. The ramp factor $\alpha(e)$ starts at the CGA start epoch and reaches 1 over 40 epochs, reducing early instability when the backbone mask predictor is still warming up.

\begin{algorithm}[t]
\caption{Training-time CGA-v2 target construction}
\label{alg:cga}
\begin{algorithmic}[1]
\REQUIRE Binary mask $y$, radii $r_c,r_b,r_p$, max area $A_{\max}$
\STATE Find eight-connected foreground components $\mathcal{C}(y)$.
\STATE Initialize $t_c,t_b,t_s,t_p$ as zero maps.
\FOR{each component $C \in \mathcal{C}(y)$}
  \STATE Mark the rounded component centroid in $t_c$.
  \STATE Set $t_s(p)=\min(1, |C|/A_{\max})$ for all $p \in C$.
\ENDFOR
\STATE Dilate $t_c$ with radius $r_c$.
\STATE Set $t_b=\mathrm{dilate}(y,r_b)-\mathrm{erode}(y,r_b)$.
\STATE Set $t_p=\mathrm{dilate}(t_c,r_p)$.
\RETURN $t_c,t_b,t_s,t_p$
\end{algorithmic}
\end{algorithm}

\subsection{Inference}

At test time, CGA-v2 follows the same inference path as MSHNetOHEM: the model returns the final mask logit, the probability map is obtained by sigmoid, and the binary mask is produced by threshold 0.5. Auxiliary logits are not used for prediction, post-processing, confidence calibration, or failure filtering. This restriction is central to the target-preserving claim: any final result must come from the same final-logit evaluation path as the baseline.

\section{Experiments}

\subsection{Evidence Status}

The current repository contains a complete current-repository MSHNetOHEM seed42 epoch400 baseline and incomplete MSHNetCGA seed42 training logs. No current-repository MSHNetCGA epoch400 checkpoint or evaluation summary is present yet. Therefore, this draft reports verified baseline values and leaves CGA, paired deltas, multi-seed statistics, and ablations as TBD. This is deliberate: historical results from other repositories and fallback/mock regularizers are not admissible paper evidence.

Table \ref{tab:evidence-gates} summarizes the evidence gates already available in the current repository. These gates do not prove CGA effectiveness, but they determine whether a completed training run may be used in the paper. The distinction matters: a model can train successfully while still being ineligible for a paper table if it uses an unverified split, a fallback regularizer, an incomplete checkpoint, or a non-fixed threshold.

\begin{table*}[t]
\centering
\caption{Current evidence gates in the repository. A pass here only authorizes the protocol; it does not replace the missing paired CGA result.}
\label{tab:evidence-gates}
\small
\begin{tabular}{lll}
\toprule
Gate & Current status & Paper implication \\
\midrule
Repository contract & Pass & Real CGA wrapper, finite loss, explicit train/eval logits. \\
Claim guard & Pass & No current repository claim-guard finding. \\
P1 dataset preflight & Pass & NUDT-SIRST split files and masks pass integrity checks. \\
P1A HC-Val source audit & Pass & HC-Val source is accepted and has zero train overlap. \\
OHEM seed42 epoch400 eval & Present & Baseline values in Table \ref{tab:ohem-baseline} are usable. \\
CGA seed42 epoch400 eval & Missing & All CGA effectiveness claims remain TBD. \\
Seed43/44 paired evidence & Missing & No multi-seed robustness claim is allowed. \\
\bottomrule
\end{tabular}
\end{table*}

\subsection{Dataset and Splits}

Experiments target NUDT-SIRST, introduced with DNANet \citep{li2023dnanet}. The local dataset preflight passes with 663 training images and 664 test images. The frozen hard-clutter validation split contains 6 images, is a subset of the test split, and has zero overlap with the training split. The HC-Val source audit also passes, with the split source accepted before the new-repository seed42 training. The current evidence protocol therefore treats the full test split as the primary result and HC-Val as a diagnostic split for hard-clutter target preservation.

\subsection{Metrics}

We report pixel-level mIoU, normalized IoU (nIoU), precision, recall, and F1, plus object-level probability of detection (Pd), false-alarm rate (FA), FA in parts per million pixels (FA ppm), and false-positive component count. Pd is computed by matching predicted connected components to target components using centroid distance, and FA/FP components are computed from unmatched predicted components. All metrics use the fixed threshold 0.5 on the final probability map.

\subsection{Implementation Protocol}

The controlled protocol uses MSHNet as the backbone, 400 epochs, Adam with learning rate $5 \times 10^{-4}$, batch size 8, patch size 256, OHEM ratio 0.01, MSHNet warm epoch 5, CGA start epoch 1, and CGA ramp 40 epochs. Paper-evidence runs require \texttt{evidence\_mode=paper}, passed P1 dataset preflight, passed P1A HC-Val source audit, and no fallback regularizer. The current seed42 CGA log reaches epoch45 in the lowercase \texttt{mshnet\_cga} run and epoch46 in the legacy-cased \texttt{MSHNetCGA} run, so neither is a complete epoch400 paper result.

The evidence package expected for each completed run contains the training log, the epoch400 checkpoint, the full-test evaluation summary, the HC-Val evaluation summary, and the run metadata recording seed, checkpoint epoch, fixed threshold, evidence mode, preflight flags, and fallback status. The evaluator writes both split-specific summaries and a compatibility full-test summary for the test split. This makes the paired comparison check mechanical: both baseline and CGA summaries must point to epoch400 checkpoints, the same dataset, the same split, the same seed, and threshold 0.5 before any delta is computed.

The controlled protocol should not be mixed with literature-only numbers. Prior methods such as MSHNet/SLS, ISNet, and PConv+SD Loss are important closest threats, but their reported literature numbers are not same-repository paired results under this frozen evaluation path. They belong in Related Work and threat analysis unless reproduced under the same data registry, checkpoint, threshold, and metric code.

\begin{table}[t]
\centering
\caption{Verified MSHNetOHEM seed42 epoch400 baseline on NUDT-SIRST in the current repository. All values use threshold 0.5 on the final logit.}
\label{tab:ohem-baseline}
\small
\begin{tabular}{lrrrr}
\toprule
Split & mIoU & Precision & Pd & FA ppm \\
\midrule
Full test & 0.9240 & 0.9692 & 0.9810 & 6.8711 \\
HC-Val & 0.8556 & 0.9275 & 0.8333 & 22.8882 \\
\bottomrule
\end{tabular}
\end{table}

\begin{table*}[t]
\centering
\caption{Planned paired seed42 comparison. CGA entries must be filled only from the current-repository MSHNetCGA seed42 epoch400 checkpoint and fixed-threshold evaluation.}
\label{tab:paired-main}
\small
\begin{tabular}{llrrrrrr}
\toprule
Split & Method & mIoU & nIoU & Precision & Recall & Pd & FA ppm \\
\midrule
Full test & MSHNetOHEM & 0.9240 & 0.9302 & 0.9692 & 0.9519 & 0.9810 & 6.8711 \\
Full test & MSHNetCGA & TBD & TBD & TBD & TBD & TBD & TBD \\
Full test & Delta & TBD & TBD & TBD & TBD & TBD & TBD \\
\midrule
HC-Val & MSHNetOHEM & 0.8556 & 0.7668 & 0.9275 & 0.9169 & 0.8333 & 22.8882 \\
HC-Val & MSHNetCGA & TBD & TBD & TBD & TBD & TBD & TBD \\
HC-Val & Delta & TBD & TBD & TBD & TBD & TBD & TBD \\
\bottomrule
\end{tabular}
\end{table*}

\subsection{Decision Rule for Completing Evidence}

The next evidence gate is the paired seed42 comparison between current-repository MSHNetOHEM and MSHNetCGA at epoch400. The predeclared Full-split gate is: mIoU delta at least +0.020, precision delta at least +0.010, and FA no worse than the baseline. HC-Val is reported as a hard-clutter diagnostic and should not be converted into an all-seed FA-decrease requirement. If seed42 fails the primary Full-split gate, the method should not proceed as an AAAI main-method route without an implementation audit.

If seed42 passes the primary gate, the next step is paired seed43/44 reproduction under the same controlled protocol. The paper should then report per-seed values, mean deltas, and the number of seeds with non-negative mIoU deltas. The multi-seed decision rule should remain weaker for HC-Val false alarms than for Full-split performance because HC-Val contains only six hard-clutter images and is meant to diagnose tradeoffs. A fair negative HC-Val row is still useful evidence if it reveals that CGA improves target preservation while increasing false alarms in a specific clutter regime.

\subsection{Ablation and Failure Analysis Plan}

After the paired seed42 result passes the primary gate, ablations should isolate the four geometry terms and the ramp schedule. Table \ref{tab:ablation-plan} records the planned ablation structure. Each row should be evaluated through the same final-logit inference path and fixed threshold; no row should use auxiliary outputs at test time.

\begin{table}[t]
\centering
\caption{Ablation table scaffold. Values are intentionally TBD until full CGA evidence exists.}
\label{tab:ablation-plan}
\small
\begin{tabular}{lrrr}
\toprule
Variant & Full mIoU & HC-Val Pd & FA ppm \\
\midrule
Full CGA & TBD & TBD & TBD \\
No center target & TBD & TBD & TBD \\
No boundary target & TBD & TBD & TBD \\
No scale target & TBD & TBD & TBD \\
No peak target & TBD & TBD & TBD \\
No ramp & TBD & TBD & TBD \\
\bottomrule
\end{tabular}
\end{table}

Failure analysis should compare matched and unmatched components on the HC-Val images. The main questions are whether CGA preserves faint target components, whether boundary supervision reduces fragmented false positives, and whether scale supervision creates a precision/FA tradeoff on hard clutter. The known paper-writing boundary is that not every seed should be required to reduce HC-Val FA; a seed may improve mIoU or Pd while worsening precision or false alarms, and such cases must be reported instead of hidden.

\subsection{Threats to Validity}

The first threat is evidence incompleteness. A regularizer paper is not validated by a clean implementation contract alone; the decisive result is whether the final mask predictor improves under paired evaluation. The current draft therefore cannot claim superiority over MSHNetOHEM or any literature method.

The second threat is split narrowness. NUDT-SIRST is the only dataset with passed preflight and HC-Val source audit in the current evidence package. NUAA-SIRST and IRSTD-1K are present as registry targets but cannot support claims until their own integrity checks, training runs, and evaluation summaries exist.

The third threat is attribution. Even after a positive full CGA result, mixed ablation outcomes are possible because center, boundary, scale, and peak targets act on a shared feature. The paper should avoid claiming strict necessity of every auxiliary target unless ablations support that statement across the fixed protocol.

The fourth threat is metric tradeoff. IRSTD performance involves pixel overlap, component detection, and false alarms. An improvement in mIoU or Pd can coexist with a worse FA or precision value on hard clutter. The final paper should report this explicitly instead of selecting a single favorable metric.

\section{Limitations and Reproducibility}

The current draft is not submission-ready because the central CGA evidence is incomplete. Its verified claims are limited to the implementation contract, passed repository/data gates, and the MSHNetOHEM seed42 baseline in Table \ref{tab:ohem-baseline}. The CGA improvement claim, multi-seed robustness, ablation attribution, and hard-clutter conclusion all require the current-repository MSHNetCGA epoch400 evaluations.

Reproducibility is part of the method boundary rather than an afterthought. Paper evidence requires the real \texttt{cga\_center\_logit}, \texttt{cga\_boundary\_logit}, \texttt{cga\_scale\_logit}, and \texttt{cga\_peak\_logit} heads; fallback regularizers are smoke-test-only and cannot support paper tables. Dataset integrity is guarded by split hashes, zero train/test overlap, zero train/HC-Val overlap, and accepted HC-Val source provenance. Evaluation uses the final logit only and records the fixed threshold.

The work uses public infrared imagery and binary target masks. It does not introduce human-subject data. The application area can include surveillance and defense-adjacent sensing, so deployment claims should remain bounded to benchmark detection behavior and should not imply operational reliability without separate validation.

\section{Conclusion}

CGA-v2 frames infrared small target detection as a target-preserving regularization problem: component-level geometry can supervise training while leaving the final mask inference path unchanged. The method attaches center, boundary, scale, and peak auxiliary heads to an explicitly declared MSHNet feature and combines their losses with the baseline OHEM and SoftIoU objective. The current repository already supports the fail-closed implementation and data gates needed for paper evidence, but the decisive paired CGA result is still missing. The next manuscript revision should replace the TBD entries with current-repository epoch400 CGA evaluations and then decide whether the evidence supports an AAAI main claim.

\bibliography{references}

\end{document}
